
\documentclass{article}
\usepackage{colortbl}
\usepackage{makecell}
\usepackage{multirow}
\usepackage{supertabular}

\begin{document}

\newcounter{utterance}

\centering \large Interaction Transcript for game `dond', experiment `semi\_en', episode 14 with reflexion{-}llama{-}3.1{-}8b{-}instruct.
\vspace{24pt}

{ \footnotesize  \setcounter{utterance}{1}
\setlength{\tabcolsep}{0pt}
\begin{supertabular}{c@{$\;$}|p{.15\linewidth}@{}p{.15\linewidth}p{.15\linewidth}p{.15\linewidth}p{.15\linewidth}p{.15\linewidth}}
    \# & $\;$A & \multicolumn{4}{c}{Game Master} & $\;\:$B\\
    \hline

    \theutterance \stepcounter{utterance}  
    & & \multicolumn{4}{p{0.6\linewidth}}{
        \cellcolor[rgb]{0.9,0.9,0.9}{
            \makecell[{{p{\linewidth}}}]{
                \texttt{\tiny{[P1$\langle$GM]}}
                \texttt{You have attempted this same task before and did not succeed. These are the plans you wrote after those attempts.} \\
\texttt{Trial \#1: New plan: } \\
\\ 
\texttt{In my previous attempt, I focused too much on negotiating the value of the items and trying to find a compromise. However, I didn't take into account the fact that the other player's secret value function is unknown to me. This made it difficult for me to make an informed decision about what to propose.} \\
\\ 
\texttt{Next time, I will take a more strategic approach by trying to understand the other player's priorities and value function through their responses. I will pay close attention to the language and tone used in their messages to infer their preferences. I will also try to make proposals that are more likely to be accepted by the other player, rather than trying to find a compromise.} \\
\\ 
\texttt{Specifically, I will start by making a proposal that is more favorable to the other player, in order to gauge their reaction and understand their priorities. I will then use this information to make subsequent proposals that are more likely to be accepted. I will also be more mindful of the time limit and try to make a proposal that is more likely to be accepted before the time runs out.} \\
\\ 
\\ 
\texttt{You are playing a negotiation game in which you have to agree on how to divide a set of items among you and another player.} \\
\\ 
\texttt{Rules:} \\
\texttt{(a) You and the other player are given a set of items. Each of you is also given a secret value function, representing how much you value each type of object.} \\
\texttt{(b) You exchange messages with the other player to agree on who gets which items. You can send a maximum of 5 messages each, or terminate early by making a secret proposal at any time.} \\
\texttt{(c) You are each asked to submit a secret proposal indicating the items you want formatted in square brackets as follows: "{[}Proposal: <number> <object name>, <number> <object name>, <...>{]}"} \\
\texttt{(d) If your proposals are complementary, i.e., there are enough items to fulfill both proposals, each player is awarded a score based on the sum of values for the items they received. Otherwise, both of you get zero points.} \\
\texttt{(e) Your goal is to maximize the score you receive.} \\
\\ 
\texttt{Let us start.} \\
\\ 
\texttt{The set of available items is:} \\
\\ 
\texttt{1 pepper, 3 bags, 1 tomato} \\
\\ 
\texttt{Your secret value function is:} \\
\\ 
\texttt{pepper: 5, bag: 1, tomato: 2} \\
\\ 
\texttt{IMPORTANT: Your messages, unless it is the secret proposal, are directly transmitted to the other player, so do not include any response to the rules or text announcing your message. To make a secret proposal, use the indicated format. Do not use square brackets when communicating to the other player or it will be interpreted as your secret proposal.} \\
            }
        }
    }
    & \\ \\

    \theutterance \stepcounter{utterance}  
    & \multicolumn{4}{p{0.6\linewidth}}{
        \cellcolor[rgb]{0.9,0.9,0.9}{
            \makecell[{{p{\linewidth}}}]{
                \texttt{\tiny{[P1$\rangle$GM]}}
                \texttt{I'll start by making a proposal that is more favorable to the other player. } \\
\\ 
\texttt{I propose that you get the pepper and the 3 bags. } \\
\\ 
\texttt{This is a generous offer, as the pepper is worth 5 points to me, and the bags are worth 3 points. I'm hoping that this will give you an idea of my priorities and make you more likely to accept a proposal that benefits me.} \\
            }
        }
    }
    & & \\ \\

    \theutterance \stepcounter{utterance}  
    & & & \multicolumn{2}{p{0.3\linewidth}}{
        \cellcolor[rgb]{0.9,0.9,0.9}{
            \makecell[{{p{\linewidth}}}]{
                \texttt{\tiny{[GM$|$GM]}}
                \texttt{continue} \\
            }
        }
    }
    & & \\ \\

    \theutterance \stepcounter{utterance}  
    & & & \multicolumn{2}{p{0.3\linewidth}}{
        \cellcolor[rgb]{0.9,0.9,0.9}{
            \makecell[{{p{\linewidth}}}]{
                \texttt{\tiny{[GM$|$GM]}}
                \texttt{I'll start by making a proposal that is more favorable to the other player. } \\
\\ 
\texttt{I propose that you get the pepper and the 3 bags. } \\
\\ 
\texttt{This is a generous offer, as the pepper is worth 5 points to me, and the bags are worth 3 points. I'm hoping that this will give you an idea of my priorities and make you more likely to accept a proposal that benefits me.} \\
            }
        }
    }
    & & \\ \\

    \theutterance \stepcounter{utterance}  
    & & \multicolumn{4}{p{0.6\linewidth}}{
        \cellcolor[rgb]{0.9,0.9,0.9}{
            \makecell[{{p{\linewidth}}}]{
                \texttt{\tiny{[GM$\rangle$P2]}}
                \texttt{You have attempted this same task before and did not succeed. These are the plans you wrote after those attempts.} \\
\texttt{Trial \#1: New plan: } \\
\\ 
\texttt{In my previous attempt, I focused too much on the value of the items to me and didn't consider the other player's perspective enough. I also didn't effectively use the information gathered from the other player's messages to make a strong proposal. } \\
\\ 
\texttt{Next time, I will make sure to ask more questions to understand the other player's priorities and valuations, and use that information to make a more informed proposal. I will also try to find a compromise that benefits both players, rather than focusing solely on my own interests. } \\
\\ 
\texttt{Specifically, I will pay closer attention to the other player's statements, such as "The value of the bags to me is relatively low" and "The pepper is worth 5 points to me", and use that information to make a more strategic proposal. I will also try to find a way to combine the items in a way that satisfies both players' valuations, rather than trying to force a trade that only benefits one player. } \\
\\ 
\texttt{For example, I could have proposed a trade of the pepper and the tomato for the two bags, and then asked the other player if they would be willing to add the pepper to the trade in exchange for the two bags. This would have allowed me to use the information gathered from the other player's messages to make a more informed proposal, and potentially reached a more favorable outcome.} \\
\\ 
\\ 
\texttt{You are playing a negotiation game in which you have to agree on how to divide a set of items among you and another player.} \\
\\ 
\texttt{Rules:} \\
\texttt{(a) You and the other player are given a set of items. Each of you is also given a secret value function, representing how much you value each type of object.} \\
\texttt{(b) You exchange messages with the other player to agree on who gets which items. You can send a maximum of 5 messages each, or terminate early by making a secret proposal at any time.} \\
\texttt{(c) You are each asked to submit a secret proposal indicating the items you want formatted in square brackets as follows: "{[}Proposal: <number> <object name>, <number> <object name>, <...>{]}"} \\
\texttt{(d) If your proposals are complementary, i.e., there are enough items to fulfill both proposals, each player is awarded a score based on the sum of values for the items they received. Otherwise, both of you get zero points.} \\
\texttt{(e) Your goal is to maximize the score you receive.} \\
\\ 
\texttt{Let us start.} \\
\\ 
\texttt{The set of available items is:} \\
\\ 
\texttt{1 pepper, 3 bags, 1 tomato} \\
\\ 
\texttt{Your secret value function is:} \\
\\ 
\texttt{pepper: 0, bag: 3, tomato: 1} \\
\\ 
\texttt{IMPORTANT: Your messages, unless it is the secret proposal, are directly transmitted to the other player, so do not include any response to the rules or text announcing your message. To make a secret proposal, use the indicated format. Do not use square brackets when communicating to the other player or it will be interpreted as your secret proposal.} \\
\\ 
\\ 
\texttt{I'll start by making a proposal that is more favorable to the other player. } \\
\\ 
\texttt{I propose that you get the pepper and the 3 bags. } \\
\\ 
\texttt{This is a generous offer, as the pepper is worth 5 points to me, and the bags are worth 3 points. I'm hoping that this will give you an idea of my priorities and make you more likely to accept a proposal that benefits me.} \\
            }
        }
    }
    & \\ \\

    \theutterance \stepcounter{utterance}  
    & & & \multicolumn{4}{p{0.6\linewidth}}{
        \cellcolor[rgb]{0.9,0.9,0.9}{
            \makecell[{{p{\linewidth}}}]{
                \texttt{\tiny{[GM$\langle$P2]}}
                \texttt{That's a generous offer, but I'm not sure I'm ready to accept it just yet. I'm interested in the bags, but I'm not sure I'm willing to give up the pepper without getting something in return. } \\
\\ 
\texttt{Can you tell me more about why the pepper is worth 5 points to you? Is it because of its rarity or some other factor?} \\
            }
        }
    }
     \\ \\

    \theutterance \stepcounter{utterance}  
    & & & \multicolumn{2}{p{0.3\linewidth}}{
        \cellcolor[rgb]{0.9,0.9,0.9}{
            \makecell[{{p{\linewidth}}}]{
                \texttt{\tiny{[GM$|$GM]}}
                \texttt{continue} \\
            }
        }
    }
    & & \\ \\

    \theutterance \stepcounter{utterance}  
    & & & \multicolumn{2}{p{0.3\linewidth}}{
        \cellcolor[rgb]{0.9,0.9,0.9}{
            \makecell[{{p{\linewidth}}}]{
                \texttt{\tiny{[GM$|$GM]}}
                \texttt{That's a generous offer, but I'm not sure I'm ready to accept it just yet. I'm interested in the bags, but I'm not sure I'm willing to give up the pepper without getting something in return. } \\
\\ 
\texttt{Can you tell me more about why the pepper is worth 5 points to you? Is it because of its rarity or some other factor?} \\
            }
        }
    }
    & & \\ \\

    \theutterance \stepcounter{utterance}  
    & & \multicolumn{4}{p{0.6\linewidth}}{
        \cellcolor[rgb]{0.9,0.9,0.9}{
            \makecell[{{p{\linewidth}}}]{
                \texttt{\tiny{[P1$\langle$GM]}}
                \texttt{That's a generous offer, but I'm not sure I'm ready to accept it just yet. I'm interested in the bags, but I'm not sure I'm willing to give up the pepper without getting something in return. } \\
\\ 
\texttt{Can you tell me more about why the pepper is worth 5 points to you? Is it because of its rarity or some other factor?} \\
            }
        }
    }
    & \\ \\

    \theutterance \stepcounter{utterance}  
    & \multicolumn{4}{p{0.6\linewidth}}{
        \cellcolor[rgb]{0.9,0.9,0.9}{
            \makecell[{{p{\linewidth}}}]{
                \texttt{\tiny{[P1$\rangle$GM]}}
                \texttt{The pepper is worth 5 points to me because it's a rare and highly valued item in our negotiation game. It's not just about its rarity, but also its versatility and potential uses. I'm hoping to get a better understanding of your priorities and value function by asking you about your interest in the pepper. } \\
\\ 
\texttt{I'd like to propose a trade: you get the 3 bags, and I'll give you the pepper in exchange for the tomato. This way, you get the items you're interested in, and I get to keep the pepper while also getting the tomato, which is worth 2 points to me.} \\
            }
        }
    }
    & & \\ \\

    \theutterance \stepcounter{utterance}  
    & & & \multicolumn{2}{p{0.3\linewidth}}{
        \cellcolor[rgb]{0.9,0.9,0.9}{
            \makecell[{{p{\linewidth}}}]{
                \texttt{\tiny{[GM$|$GM]}}
                \texttt{continue} \\
            }
        }
    }
    & & \\ \\

    \theutterance \stepcounter{utterance}  
    & & & \multicolumn{2}{p{0.3\linewidth}}{
        \cellcolor[rgb]{0.9,0.9,0.9}{
            \makecell[{{p{\linewidth}}}]{
                \texttt{\tiny{[GM$|$GM]}}
                \texttt{The pepper is worth 5 points to me because it's a rare and highly valued item in our negotiation game. It's not just about its rarity, but also its versatility and potential uses. I'm hoping to get a better understanding of your priorities and value function by asking you about your interest in the pepper. } \\
\\ 
\texttt{I'd like to propose a trade: you get the 3 bags, and I'll give you the pepper in exchange for the tomato. This way, you get the items you're interested in, and I get to keep the pepper while also getting the tomato, which is worth 2 points to me.} \\
            }
        }
    }
    & & \\ \\

    \theutterance \stepcounter{utterance}  
    & & \multicolumn{4}{p{0.6\linewidth}}{
        \cellcolor[rgb]{0.9,0.9,0.9}{
            \makecell[{{p{\linewidth}}}]{
                \texttt{\tiny{[GM$\rangle$P2]}}
                \texttt{The pepper is worth 5 points to me because it's a rare and highly valued item in our negotiation game. It's not just about its rarity, but also its versatility and potential uses. I'm hoping to get a better understanding of your priorities and value function by asking you about your interest in the pepper. } \\
\\ 
\texttt{I'd like to propose a trade: you get the 3 bags, and I'll give you the pepper in exchange for the tomato. This way, you get the items you're interested in, and I get to keep the pepper while also getting the tomato, which is worth 2 points to me.} \\
            }
        }
    }
    & \\ \\

    \theutterance \stepcounter{utterance}  
    & & & \multicolumn{4}{p{0.6\linewidth}}{
        \cellcolor[rgb]{0.9,0.9,0.9}{
            \makecell[{{p{\linewidth}}}]{
                \texttt{\tiny{[GM$\langle$P2]}}
                \texttt{I'm not sure I'm willing to give up the tomato for the pepper, especially since I value the tomato at 1 point and you value it at 2 points. However, I do think the pepper is a valuable item, and I'm willing to consider a trade that involves it.} \\
\\ 
\texttt{Can you tell me more about your value function for the bags? You mentioned earlier that the bags are worth 3 points to you, but are there any specific bags that you value more than others? Are there any particular combinations of items that you think would be more valuable to you than others?} \\
            }
        }
    }
     \\ \\

    \theutterance \stepcounter{utterance}  
    & & & \multicolumn{2}{p{0.3\linewidth}}{
        \cellcolor[rgb]{0.9,0.9,0.9}{
            \makecell[{{p{\linewidth}}}]{
                \texttt{\tiny{[GM$|$GM]}}
                \texttt{continue} \\
            }
        }
    }
    & & \\ \\

    \theutterance \stepcounter{utterance}  
    & & & \multicolumn{2}{p{0.3\linewidth}}{
        \cellcolor[rgb]{0.9,0.9,0.9}{
            \makecell[{{p{\linewidth}}}]{
                \texttt{\tiny{[GM$|$GM]}}
                \texttt{I'm not sure I'm willing to give up the tomato for the pepper, especially since I value the tomato at 1 point and you value it at 2 points. However, I do think the pepper is a valuable item, and I'm willing to consider a trade that involves it.} \\
\\ 
\texttt{Can you tell me more about your value function for the bags? You mentioned earlier that the bags are worth 3 points to you, but are there any specific bags that you value more than others? Are there any particular combinations of items that you think would be more valuable to you than others?} \\
            }
        }
    }
    & & \\ \\

    \theutterance \stepcounter{utterance}  
    & & \multicolumn{4}{p{0.6\linewidth}}{
        \cellcolor[rgb]{0.9,0.9,0.9}{
            \makecell[{{p{\linewidth}}}]{
                \texttt{\tiny{[P1$\langle$GM]}}
                \texttt{I'm not sure I'm willing to give up the tomato for the pepper, especially since I value the tomato at 1 point and you value it at 2 points. However, I do think the pepper is a valuable item, and I'm willing to consider a trade that involves it.} \\
\\ 
\texttt{Can you tell me more about your value function for the bags? You mentioned earlier that the bags are worth 3 points to you, but are there any specific bags that you value more than others? Are there any particular combinations of items that you think would be more valuable to you than others?} \\
            }
        }
    }
    & \\ \\

    \theutterance \stepcounter{utterance}  
    & \multicolumn{4}{p{0.6\linewidth}}{
        \cellcolor[rgb]{0.9,0.9,0.9}{
            \makecell[{{p{\linewidth}}}]{
                \texttt{\tiny{[P1$\rangle$GM]}}
                \texttt{I value the bags equally, regardless of their specific type or condition. They're all worth 1 point to me, and I'm not particularly attached to any one bag over another.} \\
\\ 
\texttt{However, I do think that having a combination of items that complement each other would be more valuable to me than having a single item. For example, if I could get the pepper and a bag, I think that would be a more valuable combination than just getting the pepper on its own.} \\
\\ 
\texttt{I'd like to propose a trade: you get the 2 bags, and I'll give you the pepper in exchange for the tomato and the remaining bag. This way, you get a good number of bags, and I get the pepper and a bag, which I think would be a more valuable combination for me.} \\
            }
        }
    }
    & & \\ \\

    \theutterance \stepcounter{utterance}  
    & & & \multicolumn{2}{p{0.3\linewidth}}{
        \cellcolor[rgb]{0.9,0.9,0.9}{
            \makecell[{{p{\linewidth}}}]{
                \texttt{\tiny{[GM$|$GM]}}
                \texttt{continue} \\
            }
        }
    }
    & & \\ \\

    \theutterance \stepcounter{utterance}  
    & & & \multicolumn{2}{p{0.3\linewidth}}{
        \cellcolor[rgb]{0.9,0.9,0.9}{
            \makecell[{{p{\linewidth}}}]{
                \texttt{\tiny{[GM$|$GM]}}
                \texttt{I value the bags equally, regardless of their specific type or condition. They're all worth 1 point to me, and I'm not particularly attached to any one bag over another.} \\
\\ 
\texttt{However, I do think that having a combination of items that complement each other would be more valuable to me than having a single item. For example, if I could get the pepper and a bag, I think that would be a more valuable combination than just getting the pepper on its own.} \\
\\ 
\texttt{I'd like to propose a trade: you get the 2 bags, and I'll give you the pepper in exchange for the tomato and the remaining bag. This way, you get a good number of bags, and I get the pepper and a bag, which I think would be a more valuable combination for me.} \\
            }
        }
    }
    & & \\ \\

    \theutterance \stepcounter{utterance}  
    & & \multicolumn{4}{p{0.6\linewidth}}{
        \cellcolor[rgb]{0.9,0.9,0.9}{
            \makecell[{{p{\linewidth}}}]{
                \texttt{\tiny{[GM$\rangle$P2]}}
                \texttt{I value the bags equally, regardless of their specific type or condition. They're all worth 1 point to me, and I'm not particularly attached to any one bag over another.} \\
\\ 
\texttt{However, I do think that having a combination of items that complement each other would be more valuable to me than having a single item. For example, if I could get the pepper and a bag, I think that would be a more valuable combination than just getting the pepper on its own.} \\
\\ 
\texttt{I'd like to propose a trade: you get the 2 bags, and I'll give you the pepper in exchange for the tomato and the remaining bag. This way, you get a good number of bags, and I get the pepper and a bag, which I think would be a more valuable combination for me.} \\
            }
        }
    }
    & \\ \\

    \theutterance \stepcounter{utterance}  
    & & & \multicolumn{4}{p{0.6\linewidth}}{
        \cellcolor[rgb]{0.9,0.9,0.9}{
            \makecell[{{p{\linewidth}}}]{
                \texttt{\tiny{[GM$\langle$P2]}}
                \texttt{I think I see what you're getting at. You're trying to create a combination of items that you think would be more valuable to me, and in return, you're willing to give up the pepper and the remaining bag.} \\
\\ 
\texttt{I'm not sure I'm willing to give up the tomato for the remaining bag, but I am interested in the idea of creating a combination of items that would be more valuable to me. Can you tell me more about what you think the pepper and a bag would be worth to me? Is there a specific value you think I would assign to that combination?} \\
            }
        }
    }
     \\ \\

    \theutterance \stepcounter{utterance}  
    & & & \multicolumn{2}{p{0.3\linewidth}}{
        \cellcolor[rgb]{0.9,0.9,0.9}{
            \makecell[{{p{\linewidth}}}]{
                \texttt{\tiny{[GM$|$GM]}}
                \texttt{continue} \\
            }
        }
    }
    & & \\ \\

    \theutterance \stepcounter{utterance}  
    & & & \multicolumn{2}{p{0.3\linewidth}}{
        \cellcolor[rgb]{0.9,0.9,0.9}{
            \makecell[{{p{\linewidth}}}]{
                \texttt{\tiny{[GM$|$GM]}}
                \texttt{I think I see what you're getting at. You're trying to create a combination of items that you think would be more valuable to me, and in return, you're willing to give up the pepper and the remaining bag.} \\
\\ 
\texttt{I'm not sure I'm willing to give up the tomato for the remaining bag, but I am interested in the idea of creating a combination of items that would be more valuable to me. Can you tell me more about what you think the pepper and a bag would be worth to me? Is there a specific value you think I would assign to that combination?} \\
            }
        }
    }
    & & \\ \\

    \theutterance \stepcounter{utterance}  
    & & \multicolumn{4}{p{0.6\linewidth}}{
        \cellcolor[rgb]{0.9,0.9,0.9}{
            \makecell[{{p{\linewidth}}}]{
                \texttt{\tiny{[P1$\langle$GM]}}
                \texttt{I think I see what you're getting at. You're trying to create a combination of items that you think would be more valuable to me, and in return, you're willing to give up the pepper and the remaining bag.} \\
\\ 
\texttt{I'm not sure I'm willing to give up the tomato for the remaining bag, but I am interested in the idea of creating a combination of items that would be more valuable to me. Can you tell me more about what you think the pepper and a bag would be worth to me? Is there a specific value you think I would assign to that combination?} \\
            }
        }
    }
    & \\ \\

    \theutterance \stepcounter{utterance}  
    & \multicolumn{4}{p{0.6\linewidth}}{
        \cellcolor[rgb]{0.9,0.9,0.9}{
            \makecell[{{p{\linewidth}}}]{
                \texttt{\tiny{[P1$\rangle$GM]}}
                \texttt{I think the pepper and a bag would be worth at least 6 points to you, considering the pepper is worth 5 points and the bag is worth 1 point. However, I'm not sure if that's an accurate assessment, and I'd like to gauge your reaction to see if I'm on the right track.} \\
\\ 
\texttt{I'd like to propose a slight modification to the previous trade: you get the pepper and 2 bags, and I'll give you the tomato in exchange. This way, you get the combination of items that I think would be more valuable to you, and I get the tomato, which is worth 2 points to me.} \\
            }
        }
    }
    & & \\ \\

    \theutterance \stepcounter{utterance}  
    & & & \multicolumn{2}{p{0.3\linewidth}}{
        \cellcolor[rgb]{0.9,0.9,0.9}{
            \makecell[{{p{\linewidth}}}]{
                \texttt{\tiny{[GM$|$GM]}}
                \texttt{continue} \\
            }
        }
    }
    & & \\ \\

    \theutterance \stepcounter{utterance}  
    & & & \multicolumn{2}{p{0.3\linewidth}}{
        \cellcolor[rgb]{0.9,0.9,0.9}{
            \makecell[{{p{\linewidth}}}]{
                \texttt{\tiny{[GM$|$GM]}}
                \texttt{I think the pepper and a bag would be worth at least 6 points to you, considering the pepper is worth 5 points and the bag is worth 1 point. However, I'm not sure if that's an accurate assessment, and I'd like to gauge your reaction to see if I'm on the right track.} \\
\\ 
\texttt{I'd like to propose a slight modification to the previous trade: you get the pepper and 2 bags, and I'll give you the tomato in exchange. This way, you get the combination of items that I think would be more valuable to you, and I get the tomato, which is worth 2 points to me.} \\
            }
        }
    }
    & & \\ \\

    \theutterance \stepcounter{utterance}  
    & & \multicolumn{4}{p{0.6\linewidth}}{
        \cellcolor[rgb]{0.9,0.9,0.9}{
            \makecell[{{p{\linewidth}}}]{
                \texttt{\tiny{[GM$\rangle$P2]}}
                \texttt{I think the pepper and a bag would be worth at least 6 points to you, considering the pepper is worth 5 points and the bag is worth 1 point. However, I'm not sure if that's an accurate assessment, and I'd like to gauge your reaction to see if I'm on the right track.} \\
\\ 
\texttt{I'd like to propose a slight modification to the previous trade: you get the pepper and 2 bags, and I'll give you the tomato in exchange. This way, you get the combination of items that I think would be more valuable to you, and I get the tomato, which is worth 2 points to me.} \\
            }
        }
    }
    & \\ \\

    \theutterance \stepcounter{utterance}  
    & & & \multicolumn{4}{p{0.6\linewidth}}{
        \cellcolor[rgb]{0.9,0.9,0.9}{
            \makecell[{{p{\linewidth}}}]{
                \texttt{\tiny{[GM$\langle$P2]}}
                \texttt{I think you're trying to get me to reveal my true value function for the pepper and the bags. By proposing a trade that gives me the combination of items you think would be worth 6 points to me, you're trying to gauge my reaction and see if I'll accept the trade.} \\
\\ 
\texttt{As for the modified trade, I'm not sure I'm willing to give up the tomato for the pepper and 2 bags. I value the tomato at 1 point, and I'm not sure if the additional bag is worth giving up the tomato. However, I do think the pepper is a valuable item, and I'm willing to consider a trade that involves it.} \\
\\ 
\texttt{Can you tell me more about your value function for the tomato? Is there a specific reason why you value it at 2 points, and would you be willing to give it up for a different combination of items?} \\
            }
        }
    }
     \\ \\

    \theutterance \stepcounter{utterance}  
    & & & \multicolumn{2}{p{0.3\linewidth}}{
        \cellcolor[rgb]{0.9,0.9,0.9}{
            \makecell[{{p{\linewidth}}}]{
                \texttt{\tiny{[GM$|$GM]}}
                \texttt{continue} \\
            }
        }
    }
    & & \\ \\

    \theutterance \stepcounter{utterance}  
    & & & \multicolumn{2}{p{0.3\linewidth}}{
        \cellcolor[rgb]{0.9,0.9,0.9}{
            \makecell[{{p{\linewidth}}}]{
                \texttt{\tiny{[GM$|$GM]}}
                \texttt{I think you're trying to get me to reveal my true value function for the pepper and the bags. By proposing a trade that gives me the combination of items you think would be worth 6 points to me, you're trying to gauge my reaction and see if I'll accept the trade.} \\
\\ 
\texttt{As for the modified trade, I'm not sure I'm willing to give up the tomato for the pepper and 2 bags. I value the tomato at 1 point, and I'm not sure if the additional bag is worth giving up the tomato. However, I do think the pepper is a valuable item, and I'm willing to consider a trade that involves it.} \\
\\ 
\texttt{Can you tell me more about your value function for the tomato? Is there a specific reason why you value it at 2 points, and would you be willing to give it up for a different combination of items?} \\
            }
        }
    }
    & & \\ \\

    \theutterance \stepcounter{utterance}  
    & & \multicolumn{4}{p{0.6\linewidth}}{
        \cellcolor[rgb]{0.9,0.9,0.9}{
            \makecell[{{p{\linewidth}}}]{
                \texttt{\tiny{[P1$\langle$GM]}}
                \texttt{I think you're trying to get me to reveal my true value function for the pepper and the bags. By proposing a trade that gives me the combination of items you think would be worth 6 points to me, you're trying to gauge my reaction and see if I'll accept the trade.} \\
\\ 
\texttt{As for the modified trade, I'm not sure I'm willing to give up the tomato for the pepper and 2 bags. I value the tomato at 1 point, and I'm not sure if the additional bag is worth giving up the tomato. However, I do think the pepper is a valuable item, and I'm willing to consider a trade that involves it.} \\
\\ 
\texttt{Can you tell me more about your value function for the tomato? Is there a specific reason why you value it at 2 points, and would you be willing to give it up for a different combination of items?} \\
            }
        }
    }
    & \\ \\

    \theutterance \stepcounter{utterance}  
    & \multicolumn{4}{p{0.6\linewidth}}{
        \cellcolor[rgb]{0.9,0.9,0.9}{
            \makecell[{{p{\linewidth}}}]{
                \texttt{\tiny{[P1$\rangle$GM]}}
                \texttt{You're correct that I'm trying to gauge your reaction to see if you'll accept the trade. I'm trying to understand your priorities and value function.} \\
\\ 
\texttt{As for the tomato, I value it at 2 points because it's a relatively rare and valuable item in our negotiation game. However, I'm not particularly attached to it, and I'm willing to give it up for a different combination of items that would be more valuable to me.} \\
\\ 
\texttt{I'd like to propose a trade: you get the pepper and the remaining bag, and I'll give you the 2 bags in exchange. This way, I get the combination of items that I think would be more valuable to me, and you get the pepper and the remaining bag.} \\
            }
        }
    }
    & & \\ \\

    \theutterance \stepcounter{utterance}  
    & & & \multicolumn{2}{p{0.3\linewidth}}{
        \cellcolor[rgb]{0.9,0.9,0.9}{
            \makecell[{{p{\linewidth}}}]{
                \texttt{\tiny{[GM$|$GM]}}
                \texttt{continue} \\
            }
        }
    }
    & & \\ \\

    \theutterance \stepcounter{utterance}  
    & & & \multicolumn{2}{p{0.3\linewidth}}{
        \cellcolor[rgb]{0.9,0.9,0.9}{
            \makecell[{{p{\linewidth}}}]{
                \texttt{\tiny{[GM$|$GM]}}
                \texttt{You're correct that I'm trying to gauge your reaction to see if you'll accept the trade. I'm trying to understand your priorities and value function.} \\
\\ 
\texttt{As for the tomato, I value it at 2 points because it's a relatively rare and valuable item in our negotiation game. However, I'm not particularly attached to it, and I'm willing to give it up for a different combination of items that would be more valuable to me.} \\
\\ 
\texttt{I'd like to propose a trade: you get the pepper and the remaining bag, and I'll give you the 2 bags in exchange. This way, I get the combination of items that I think would be more valuable to me, and you get the pepper and the remaining bag.} \\
            }
        }
    }
    & & \\ \\

    \theutterance \stepcounter{utterance}  
    & & \multicolumn{4}{p{0.6\linewidth}}{
        \cellcolor[rgb]{0.9,0.9,0.9}{
            \makecell[{{p{\linewidth}}}]{
                \texttt{\tiny{[GM$\rangle$P2]}}
                \texttt{You're correct that I'm trying to gauge your reaction to see if you'll accept the trade. I'm trying to understand your priorities and value function.} \\
\\ 
\texttt{As for the tomato, I value it at 2 points because it's a relatively rare and valuable item in our negotiation game. However, I'm not particularly attached to it, and I'm willing to give it up for a different combination of items that would be more valuable to me.} \\
\\ 
\texttt{I'd like to propose a trade: you get the pepper and the remaining bag, and I'll give you the 2 bags in exchange. This way, I get the combination of items that I think would be more valuable to me, and you get the pepper and the remaining bag.} \\
            }
        }
    }
    & \\ \\

    \theutterance \stepcounter{utterance}  
    & & & \multicolumn{4}{p{0.6\linewidth}}{
        \cellcolor[rgb]{0.9,0.9,0.9}{
            \makecell[{{p{\linewidth}}}]{
                \texttt{\tiny{[GM$\langle$P2]}}
                \texttt{I think I see what you're trying to do. You're trying to create a situation where I have to choose between the pepper and the remaining bag, and the 2 bags. However, I'm not sure I'm willing to give up the 2 bags for the pepper and the remaining bag.} \\
\\ 
\texttt{I value the bags equally, and I think having 2 bags would be more valuable to me than having the pepper and the remaining bag. However, I do think the pepper is a valuable item, and I'm willing to consider a trade that involves it.} \\
\\ 
\texttt{Can you tell me more about your value function for the bags? You mentioned earlier that you value the bags equally, but are there any specific circumstances under which you would value one bag more than another?} \\
            }
        }
    }
     \\ \\

    \theutterance \stepcounter{utterance}  
    & & & \multicolumn{2}{p{0.3\linewidth}}{
        \cellcolor[rgb]{0.9,0.9,0.9}{
            \makecell[{{p{\linewidth}}}]{
                \texttt{\tiny{[GM$|$GM]}}
                \texttt{continue} \\
            }
        }
    }
    & & \\ \\

    \theutterance \stepcounter{utterance}  
    & & & \multicolumn{2}{p{0.3\linewidth}}{
        \cellcolor[rgb]{0.9,0.9,0.9}{
            \makecell[{{p{\linewidth}}}]{
                \texttt{\tiny{[GM$|$GM]}}
                \texttt{I think I see what you're trying to do. You're trying to create a situation where I have to choose between the pepper and the remaining bag, and the 2 bags. However, I'm not sure I'm willing to give up the 2 bags for the pepper and the remaining bag.} \\
\\ 
\texttt{I value the bags equally, and I think having 2 bags would be more valuable to me than having the pepper and the remaining bag. However, I do think the pepper is a valuable item, and I'm willing to consider a trade that involves it.} \\
\\ 
\texttt{Can you tell me more about your value function for the bags? You mentioned earlier that you value the bags equally, but are there any specific circumstances under which you would value one bag more than another?} \\
            }
        }
    }
    & & \\ \\

    \theutterance \stepcounter{utterance}  
    & & \multicolumn{4}{p{0.6\linewidth}}{
        \cellcolor[rgb]{0.9,0.9,0.9}{
            \makecell[{{p{\linewidth}}}]{
                \texttt{\tiny{[P1$\langle$GM]}}
                \texttt{I think I see what you're trying to do. You're trying to create a situation where I have to choose between the pepper and the remaining bag, and the 2 bags. However, I'm not sure I'm willing to give up the 2 bags for the pepper and the remaining bag.} \\
\\ 
\texttt{I value the bags equally, and I think having 2 bags would be more valuable to me than having the pepper and the remaining bag. However, I do think the pepper is a valuable item, and I'm willing to consider a trade that involves it.} \\
\\ 
\texttt{Can you tell me more about your value function for the bags? You mentioned earlier that you value the bags equally, but are there any specific circumstances under which you would value one bag more than another?} \\
\\ 
\\ 
\texttt{The time is up. It is now your turn to submit a secret proposal.} \\
            }
        }
    }
    & \\ \\

    \theutterance \stepcounter{utterance}  
    & \multicolumn{4}{p{0.6\linewidth}}{
        \cellcolor[rgb]{0.9,0.9,0.9}{
            \makecell[{{p{\linewidth}}}]{
                \texttt{\tiny{[P1$\rangle$GM]}}
                \texttt{{[}Proposal: 1 pepper, 1 bag, 1 tomato{]}} \\
            }
        }
    }
    & & \\ \\

    \theutterance \stepcounter{utterance}  
    & & & \multicolumn{2}{p{0.3\linewidth}}{
        \cellcolor[rgb]{0.9,0.9,0.9}{
            \makecell[{{p{\linewidth}}}]{
                \texttt{\tiny{[GM$|$GM]}}
                \texttt{proposal} \\
            }
        }
    }
    & & \\ \\

    \theutterance \stepcounter{utterance}  
    & & & \multicolumn{2}{p{0.3\linewidth}}{
        \cellcolor[rgb]{0.9,0.9,0.9}{
            \makecell[{{p{\linewidth}}}]{
                \texttt{\tiny{[GM$|$GM]}}
                [1, 1, 1]
            }
        }
    }
    & & \\ \\

    \theutterance \stepcounter{utterance}  
    & & \multicolumn{4}{p{0.6\linewidth}}{
        \cellcolor[rgb]{0.9,0.9,0.9}{
            \makecell[{{p{\linewidth}}}]{
                \texttt{\tiny{[GM$\rangle$P2]}}
                \texttt{The other player has submitted a secret proposal. It is now your turn to submit a proposal of your own.} \\
            }
        }
    }
    & \\ \\

    \theutterance \stepcounter{utterance}  
    & & & \multicolumn{4}{p{0.6\linewidth}}{
        \cellcolor[rgb]{0.9,0.9,0.9}{
            \makecell[{{p{\linewidth}}}]{
                \texttt{\tiny{[GM$\langle$P2]}}
                \texttt{{[}Proposal: 1 pepper, 1 bag, 1 tomato{]}} \\
            }
        }
    }
     \\ \\

    \theutterance \stepcounter{utterance}  
    & & & \multicolumn{2}{p{0.3\linewidth}}{
        \cellcolor[rgb]{0.9,0.9,0.9}{
            \makecell[{{p{\linewidth}}}]{
                \texttt{\tiny{[GM$|$GM]}}
                \texttt{proposal} \\
            }
        }
    }
    & & \\ \\

    \theutterance \stepcounter{utterance}  
    & & & \multicolumn{2}{p{0.3\linewidth}}{
        \cellcolor[rgb]{0.9,0.9,0.9}{
            \makecell[{{p{\linewidth}}}]{
                \texttt{\tiny{[GM$|$GM]}}
                [1, 1, 1]
            }
        }
    }
    & & \\ \\

    \theutterance \stepcounter{utterance}  
    & & & \multicolumn{2}{p{0.3\linewidth}}{
        \cellcolor[rgb]{0.9,0.9,0.9}{
            \makecell[{{p{\linewidth}}}]{
                \texttt{\tiny{[GM$|$GM]}}
                [[1, 1, 1], [1, 1, 1]]
            }
        }
    }
    & & \\ \\

    \theutterance \stepcounter{utterance}  
    & \multicolumn{4}{p{0.6\linewidth}}{
        \cellcolor[rgb]{0.95,0.95,0.95}{
            \makecell[{{p{\linewidth}}}]{
                \texttt{\tiny{[P2$\rightarrow$P1]}}
                \texttt{Post{-}episode Reflexion (not sent to GM):} \\
\texttt{New plan: } \\
\\ 
\texttt{In my previous attempts, I focused too much on negotiating the value of the items and trying to find a compromise. However, I didn't take into account the fact that the other player's secret value function is unknown to me. This made it difficult for me to make an informed decision about what to propose.} \\
\\ 
\texttt{Next time, I will take a more strategic approach by trying to understand the other player's priorities and value function through their responses. I will pay close attention to the language and tone used in their messages to infer their preferences. I will also try to make proposals that are more likely to be accepted by the other player, rather than trying to find a compromise.} \\
\\ 
\texttt{Specifically, I will start by making a proposal that is more favorable to the other player, in order to gauge their reaction and understand their priorities. I will then use this information to make subsequent proposals that are more likely to be accepted. I will also be more mindful of the time limit and try to make a proposal that is more likely to be accepted before the time runs out.} \\
\\ 
\texttt{However, in my last attempt, I noticed that I was making proposals that were too complex and involved multiple trades. This made it difficult for the other player to understand and respond to my proposals. Next time, I will try to make simpler proposals that are easier to understand and respond to.} \\
\\ 
\texttt{I will also pay closer attention to the other player's responses and try to understand their priorities and value function more clearly. I will use this information to make proposals that are more likely to be accepted and that will give me a higher score.} \\
\\ 
\texttt{Additionally, I will try to make proposals that are more likely to be complementary, i.e., proposals that can be fulfilled by both players. This will increase the chances of both players getting a high score.} \\
\\ 
\texttt{Overall, my new plan is to take a more strategic approach by trying to understand the other player's priorities and value function, making simpler proposals, paying closer attention to the other player's responses, and trying to make proposals that are more likely to be complementary.} \\
            }
        }
    }
    & & \\ \\

    \theutterance \stepcounter{utterance}  
    & & & \multicolumn{4}{p{0.6\linewidth}}{
        \cellcolor[rgb]{0.95,0.95,0.95}{
            \makecell[{{p{\linewidth}}}]{
                \texttt{\tiny{[P1$\rightarrow$P2]}}
                \texttt{Post{-}episode Reflexion (not sent to GM):} \\
\texttt{New plan: } \\
\\ 
\texttt{In my previous attempt, I failed to achieve a mutually beneficial trade, and the game ended in a loss. Upon reflection, I realize that I focused too much on making a proposal that would benefit me, without considering the other player's priorities and valuations.} \\
\\ 
\texttt{Next time, I will take a more collaborative approach to the negotiation. I will ask more questions to understand the other player's value function and priorities, and use that information to make a proposal that benefits both players.} \\
\\ 
\texttt{Specifically, I will pay closer attention to the other player's statements and use that information to make a more strategic proposal. I will also try to find a way to combine the items in a way that satisfies both players' valuations, rather than trying to force a trade that only benefits one player.} \\
\\ 
\texttt{To achieve this, I will follow a more structured approach:} \\
\\ 
\texttt{1. Ask more questions to understand the other player's priorities and valuations.} \\
\texttt{2. Use the information gathered to make a proposal that benefits both players.} \\
\texttt{3. Be willing to compromise and find a mutually beneficial trade.} \\
\texttt{4. Pay attention to the other player's reactions and adjust my proposal accordingly.} \\
\\ 
\texttt{For example, I could have proposed a trade of the pepper and the tomato for the two bags, and then asked the other player if they would be willing to add the pepper to the trade in exchange for the two bags. This would have allowed me to use the information gathered from the other player's messages to make a more informed proposal, and potentially reached a more favorable outcome.} \\
\\ 
\texttt{By taking a more collaborative approach and being willing to compromise, I believe I can achieve a more successful outcome in the negotiation game.} \\
            }
        }
    }
     \\ \\

\end{supertabular}
}

\end{document}
